一个急诊科同时盯着三类患者的二十四小时恶化风险。第一类的结局被完整记录下来，是否恶化写在表里，Logistic 回归直接拟合风险分数就够用。第二类只留下了分诊等级，而等级并不等于实情——轻症被误标为重症的比例谁也说不清，于是``患者究竟属于哪一类''本身成了未知量。第三类还要每小时测一次生命体征，体征的变化取决于一个看不见的、随时间演进的状态。三类问题的数据结构不同，可用的模型却可以出自同一个模板。

这个模板就是生成式建模。判别式模型学的是从输入到输出的那条映射，它只回答被问到的那个问题，省力而稳健。生成式模型写下的是完整的联合分布：类别怎样出现、隐状态怎样转移、观测怎样从状态里产生。写下联合分布之后，采样、补缺失、给不确定性都变成同一个模型上的不同查询，标签缺失时也还有话可说。代价同样明确——你要为所有变量作出分布假设，假设越多，能出错的地方越多。

\begin{center}
\begin{tikzpicture}[>=stealth]

% (a) 标签可见
\begin{scope}
\fill[black!18] (0.9,2.2) circle (0.32);
\draw[black!70] (0.9,2.2) circle (0.32);
\node[font=\scriptsize] at (0.9,2.2) {\(y\)};
\fill[black!18] (0.9,0.9) circle (0.32);
\draw[black!70] (0.9,0.9) circle (0.32);
\node[font=\scriptsize] at (0.9,0.9) {\(x\)};
\draw[black!70,->] (0.9,1.86) -- (0.9,1.24);
\node[font=\scriptsize,anchor=north] at (0.9,0.35) {标签可见};
\end{scope}

% (b) 标签隐藏
\begin{scope}[shift={(3.4,0)}]
\draw[black!70] (0.9,2.2) circle (0.32);
\node[font=\scriptsize] at (0.9,2.2) {\(z\)};
\fill[black!18] (0.9,0.9) circle (0.32);
\draw[black!70] (0.9,0.9) circle (0.32);
\node[font=\scriptsize] at (0.9,0.9) {\(x\)};
\draw[black!70,->] (0.9,1.86) -- (0.9,1.24);
\node[font=\scriptsize,anchor=north] at (0.9,0.35) {标签隐藏};
\end{scope}

% (c) 隐变量成链
\begin{scope}[shift={(6.4,0)}]
\foreach \a/\t in {0.9/1, 2.4/2, 3.9/3}{
  \draw[black!70] (\a,2.2) circle (0.32);
  \node[font=\scriptsize] at (\a,2.2) {\(z_{\t}\)};
  \fill[black!18] (\a,0.9) circle (0.32);
  \draw[black!70] (\a,0.9) circle (0.32);
  \node[font=\scriptsize] at (\a,0.9) {\(x_{\t}\)};
  \draw[black!70,->] (\a,1.86) -- (\a,1.24);
}
\draw[black!70,->] (1.24,2.2) -- (2.06,2.2);
\draw[black!70,->] (2.74,2.2) -- (3.56,2.2);
\node[font=\scriptsize,anchor=north] at (2.4,0.35) {隐变量沿时间成链};
\end{scope}

\end{tikzpicture}
\end{center}

\emph{灰底表示能看见，白底表示看不见；这一章沿着从左到右的方向，逐步把信息拿走。}

图里那条从左到右的路线就是本章的顺序。\hyperref[sec:13-Machine-Learning/06-Probabilistic-Models/01]{``生成式分类从 Naive Bayes 到 GDA''}处理最左边的情形：标签可见，先描述每一类怎样产生特征，再用 Bayes 公式反推类别；条件独立与 Gaussian 假设各自能买到什么、各自在哪里失效，在这一节说清。\hyperref[sec:13-Machine-Learning/06-Probabilistic-Models/03]{``Gaussian 混合模型让类别成为隐变量''}把标签遮住，用后验责任保留软归属。\hyperref[sec:13-Machine-Learning/06-Probabilistic-Models/04]{``EM 用下界交替完成隐变量学习''}解开``参数未知''与``归属未知''的循环依赖。\hyperref[sec:13-Machine-Learning/06-Probabilistic-Models/05]{``HMM 把隐变量排列成时间链''}加进时间依赖，并用动态规划把指数条路径压成每列几个数。第一次读，这四篇连起来读完即可，它们共同回答一个问题：信息被藏起来之后，怎么还能学。

其余三篇解释``为什么会长成这样''，适合第二遍回补。\hyperref[sec:13-Machine-Learning/06-Probabilistic-Models/02]{``LDA 与收缩：判别方向的几何与高维失效''}把线性判别方向还原成白化之后的均值连线，并说明为什么在高维小样本下崩掉的总是求逆那一步。\hyperref[sec:13-Machine-Learning/06-Probabilistic-Models/06]{``指数族把归一化、矩与似然连接起来''}指出前面所有 M 步的闭式解出自同一条恒等式：对数配分函数的梯度就是充分统计量的期望。\hyperref[sec:13-Machine-Learning/06-Probabilistic-Models/07]{``共轭先验让 Bayesian 更新保持同一形式''}接着说明，先验若与似然共用同一组坐标，Bayes 更新就退化成两个计数相加。

有几条判断贯穿全章，值得先摆出来。隐变量是一件建模装置，不是被发现的实体——它的作用是把一个复杂的边缘分布写成若干简单成分的混合，混合高斯、主题模型、因子分析都在做这件事。混合模型的成分可以整体重排而似然不变，所以``第 3 个成分''在跨运行时没有意义，比较两次拟合必须先做匹配。EM 是坐标上升，E 步把下界顶到与目标相切，M 步沿这张下界爬到顶点，似然因此单调不减，但只保证收敛到驻点；那个相切的几何图景比递推公式更值得记住。至于模型选择，似然更高从来不等于模型更对：成分数的选择需要额外准则，而不同准则给出的答案常常不一致，重跑一次也未必稳定。

读每一节时，把三件事分开问：模型允许哪些分布，算法实际算出了什么，现有证据凭什么支持你给隐变量起的那个名字。计算上的便利只能回答第二问，它替代不了模型检验，更替代不了外部证据。
